\section{Future ideas for Proving the Universal Boundary Inequality}\label{app:alt-strategies}

\begin{quote}
\textit{Everything in this appendix is a research sketch, not a result.
Nothing here is proven, checked, or claimed to be correct; it is recorded
only so that future work on Proposition~\ref{prop:universal} does not
restart from nothing. Neither strategy below has been carried past an
informal whiteboard stage.}
\end{quote}

The coordinate-root compression approach outlined in the accompanying
research notes (\texttt{docs/proof-\allowbreak sketch-\allowbreak weighted-potential.md}) stalls
on a coordinate-tangling obstacle: a local compression move at position
$p$ unavoidably scrambles roots at every other position $q \neq p$, and no
argument to date bounds that scatter damage tightly enough to recover
monotonicity of $X(V') + (m+1)D(V')$. Three alternative directions are
sketched below.

\paragraph{Marginal induction with partial-layer tracking.}
This direction extends the proof skeleton already used for
Theorem~\ref{thm:defect} (proven by strong induction on $R$, partitioning
$V'$ by a coordinate and applying Lemma~\ref{lem:subadditivity}), carrying
the joint quantity $\Phi = X + (m+1)D$ instead of $D$ alone. Stripping a
vertex $v$ with $s$ roots shared with the remaining $R-1$ vertices gives
$\Delta D = s$ exactly; closing the induction would require bounding
$\Delta X$ as a function of $s$ and $m$ alone, so that
$\Delta X + (m+1)s \le \Delta \Cconst(R) + m \cdot \Delta \Eseq(R)$. A
partial result toward this bound has since been derived (not yet
mechanized in Lean): writing out the marginal change in
$\mathrm{total\_coord\_edges}$ gives $\Delta X \le (k-s)m$ unconditionally,
and separately $\Delta X = 0$ whenever $v$ is adjacent to the vertex
removed (proved directly from the one-coordinate-difference property of
adjacency, independent of $n,k,R$). A second, dimension-independent cap
$\Delta X \le 2(R-1)$ also holds, from the fact that only vertices at
Hamming distance $2$ from $v$ can contribute to $\Delta X$ via a fresh
root. Whether $\min\big((k-s)m,\,2(R-1)\big)$ always suffices to close the
induction is conjectured but unproved; only two supporting configurations
have been checked (see \texttt{docs/proof-sketch-weighted-potential.md}
for the derivations and the open reconciliation with $\Delta \Cconst(R)$).

\paragraph{Dual root-compression via averaging.}
Inspired by Pinto's resolution of the Bollob\'{a}s--Leader directed-path
conjectures on the hypercube \cite{pinto2015directed}, where two
compression operators that each individually fail to shrink a directed
boundary are shown to satisfy an averaging inequality forcing at least
one of them to be non-worsening. The single guarded-symbol compression
already refuted on $A(n,k)$ (Section~\ref{sec:nogo}) is one such
individually-failing operator; whether a dual pair of root-compression
operators exists on $A(n,k)$, respecting the injectivity constraint, for
which an analogous averaging bound holds for $X+(m+1)D$, is entirely
untried. This is flagged only as a structural analogy worth investigating
for the coordinate-tangling obstacle; no candidate operator pair has been
proposed.

\paragraph{A global Lyapunov function on fiber multisets.}
This direction proposes a potential $\Phi(V') = \sum_{p=1}^k \sum_r
f(|F_{p,r}|)$ over the multiset of fiber sizes at every coordinate, for
some convex $f$, and asks whether a single root-compression move's local
gain (from concentrating vertices into larger fibers at one coordinate)
always dominates its collateral loss (from scrambling roots at every other
coordinate). Neither the link between this $\Phi$ and $X + (m+1)D$, nor
the required domination bound, has been established; it is not known
whether a suitable $f$ exists.

\medskip
\noindent Of the three, the marginal-induction direction is judged most
advanced in practice, since its early steps reuse machinery already
mechanized in Lean (Lemma~\ref{lem:subadditivity} and
Theorem~\ref{thm:defect}'s induction) and its open step now has a partial
derivation rather than being untouched. The dual-compression direction is
flagged as a promising untried structural analogy, precisely because it
targets the coordinate-tangling obstacle that stalled the original
single-operator compression, but nothing on $A(n,k)$ has been attempted
for it yet. The Lyapunov direction remains a from-scratch invariant
search with no candidate potential shown to work. These are research
opinions, not mathematical claims, and should not be read as evidence
toward Proposition~\ref{prop:universal} itself.
